\documentclass[conference]{IEEEtran}
\IEEEoverridecommandlockouts
% The preceding line is only needed to identify funding in the first footnote. If that is unneeded, please comment it out.
\usepackage{cite}
\usepackage{amsmath,amssymb,amsfonts}
\usepackage{algorithmic}
\usepackage{graphicx}
\usepackage{textcomp}
\usepackage{xcolor}
\usepackage{booktabs}

\usepackage{ctex} %%编译中文 、XeLaTeX
\usepackage{unicode-math}
\usepackage[ruled,vlined,linesnumbered]{algorithm2e}
\usepackage{amsmath}
\usepackage{amssymb}
\renewcommand{\figurename}{Figure}
\renewcommand{\tablename}{Table}
\renewcommand{\abstractname}{Abstract}
\def\BibTeX{{\rm B\kern-.05em{\sc i\kern-.025em b}\kern-.08em
    T\kern-.1667em\lower.7ex\hbox{E}\kern-.125emX}}
\renewcommand{\refname}{References}
\begin{document}

\title{Hash–Based Approximate Annulus Query for  Kernel Density Estimation in High-Dimensional Consumer IoT Data Modeling
}

\author{\IEEEauthorblockN{1\textsuperscript{st} Given Name Surname}
\IEEEauthorblockA{\textit{dept. name of organization (of Aff.)} \\
\textit{name of organization (of Aff.)}\\
City, Country \\
email address or ORCID}
\and
\IEEEauthorblockN{2\textsuperscript{nd} Given Name Surname}
\IEEEauthorblockA{\textit{dept. name of organization (of Aff.)} \\
\textit{name of organization (of Aff.)}\\
City, Country \\
email address or ORCID}

}

\maketitle

\begin{abstract}
Kernel density estimation (KDE) is widely used for high-dimensional data modeling and anomaly detection in consumer electronics, but exact KDE becomes costly when the sample size and feature dimension increase. This paper proposes an approximate annulus KDE framework based on Multi-Encoder Contrastive Hashing (MECH). MECH uses multiple encoder branches with a shared decoder to generate complementary hash tables, and combines hash-based angular retrieval, sphere-layer radial filtering, and difference-based annulus construction. To connect annulus retrieval quality with KDE accuracy, we evaluate not only point-level precision, recall, and F1-score, but also kernel-weighted recall, kernel-weighted false-positive ratio, KDE relative error, candidate expansion ratio, and online query time. Experiments on ISOLET show that after translating the normalized data away from the origin, Full MECH achieves 96.75\% recall, 96.76\% kernel-weighted recall, and 1.43\% KDE relative error in approximate annulus retrieval. The structure ablation further shows that the full combination of hash retrieval, distance-table angular filtering, sphere-layer radial filtering, and difference-based annulus construction is substantially more accurate than incomplete variants. On the UNSW-NB15 anomaly detection task, the approximate annulus KDE method achieves an F1-score of 98.1\% and provides a substantial speedup over exact annulus KDE.
\end{abstract}


\begin{IEEEkeywords}
High-dimensional Data, Kernel Density Estimation, Approximate Annulus Partition, Anomaly Detection
\end{IEEEkeywords}



\section{Introduction}
With the widespread adoption of consumer electronics such as smartphones, wearable devices, and smart home products, a large amount of high-dimensional data has been generated, including sensor signals, user interaction behaviors, and multimodal content. How to perform efficient and accurate data modeling and anomaly detection in high-dimensional spaces has become a key issue in the development of intelligent consumer electronics. KDE, as a representative non-parametric distribution modeling tool, has demonstrated strong adaptability in tasks such as anomaly detection \cite{reff1}, clustering \cite{reff2,reff3}, and pattern recognition \cite{reff4}. However, consumer electronic devices are constrained by computational power, while also requiring high real-time performance, making the high computational complexity of traditional KDE difficult to meet practical needs. Therefore, achieving efficient approximate estimation under resource constraints and low-latency requirements has become a crucial research direction.

Existing optimization methods for high-dimensional KDE can be categorized into three main types: space partition, dimensionality reduction, and nearest-neighbor search.
Space partition methods (KD-Tree \cite{reff6}, Random Tree \cite{reff7}) perform local computations by dividing the data;
Dimensionality reduction methods (PCA \cite{reff10}, manifold learning \cite{reff12}, deep autoencoders \cite{reff9}) map high-dimensional data to a lower-dimensional space;
Nearest-neighbor search methods (quantization \cite{reff15,reff16}, LSH \cite{reff17,reff18}) accelerate kernel value computation through index structures.
Although these methods significantly improve computational efficiency, it remains challenging to balance accuracy and efficiency in high-dimensional scenarios, mainly facing the following three issues:

The number of sampled points depends on the dataset size, leading to degraded performance.
In traditional KDE methods, to accurately compute the kernel density value of a query point, the algorithm typically samples from the entire dataset or partitioned subsets for approximation. However, as the dataset grows, the required number of samples increases linearly or even superlinearly.

Second, space partition tends to degenerate in high-dimensional scenarios.
Common spatial indexing and partition methods (KD-tree, random tree, random projection tree) tend to degrade into near-linear scans on sparse high-dimensional data, increasing query and maintenance overhead and making it difficult to maintain stable acceleration benefits. This degradation directly limits the overall performance of approximate annulus queries and KDE.

Third, it is difficult to effectively balance the number and quality of spatially partitioned subsets.
To ensure that sampling from the partitioned subsets can effectively estimate the kernel density values, each subset must maintain a certain level of quality.
However, ensuring high subset quality makes it difficult to guarantee sufficient partition quantity, while maintaining a large number of partitions compromises subset quality.

To address these challenges, this paper proposes a high-dimensional KDE method based on approximate annulus partition.
The proposed method leverages deep hash encoding and distance transformation to achieve efficient annulus partition and performs random sampling within each annulus. Under controllable error constraints, it significantly reduces the number of samples while maintaining high accuracy, demonstrating superior performance in anomaly detection tasks. Our contributions are as follows. 
\begin{itemize}
 \item Reducing dependence on the number of sampled points and improving the efficiency of large-scale data processing.
To address the performance degradation caused by the increase in sampling points with dataset size, this paper proposes a KDE framework based on approximate annulus sampling.
Through an annulus-level sampling strategy, the method achieves estimation accuracy close to that of exact approaches with a significantly reduced number of samples under given error constraints, thereby effectively reducing computational complexity and improving processing efficiency on large-scale datasets.

 \item 本文设计了一种“角空间到欧氏空间”的候选集构造过程。球面表将候选点限制在可行的径向层内，距离表则通过哈希距离查询候选点限制在目标角度范围内，使环半径可以在查询阶段指定。
 \item  We propose a deep hash algorithm based on high-dimensional vectors, addressing the problem of spatial partition degradation in high-dimensional scenarios,
The algorithm employs a deep neural network to perform nonlinear mapping and compression of high-dimensional features, generating compact binary hash codes that allow similar data points to fall into the same or neighboring hash buckets with higher probability in the hash space.
During querying, only a limited number of hash buckets need to be compared instead of traversing the original high-dimensional space, thereby mitigating the degradation problem of traditional KD-Trees and Random Trees on sparse high-dimensional data.


 \item  Balancing the number and quality of subsets to achieve both accuracy and efficiency, by optimizing the subset range and quantity in the design of annulus partition and sampling.
This ensures that each subset maintains a small range to guarantee kernel density estimation accuracy, while keeping the number of subsets controllable, thus achieving a balance between estimation accuracy and computational efficiency.
\end{itemize}

\section{Related Work}
In recent years, KDE has been widely applied in the field of consumer electronics, particularly in anomaly detection, wireless localization, and sensor monitoring.
\cite{shylendra2020low} utilized non-parametric KDE combined with a sliding window mechanism to model the statistical characteristics of sensor data streams, achieving unsupervised anomaly detection and ultra-low power consumption on embedded devices, making it suitable for wearable and IoT scenarios.
\cite{xu2020compressed} proposed compressed multivariate KDE for modeling WiFi fingerprint distributions, which reduces storage and computational costs while maintaining high localization accuracy on resource-constrained mobile devices.
\cite{yong2024kernel} introduced optimized fingerprint points in high-density trajectory regions, demonstrating KDE’s ability to characterize non-Gaussian and spatially heterogeneous distributions.
In addition, \cite{hu2018anomaly} combined local KDE with contextual regression to identify anomalies in the feature space, while \cite{rosenberger2022extended} proposed an Extended Kernel Density Estimation (EEM-KDE) method capable of capturing concept drift and collective anomalies, making it more suitable for real-time sensor and network monitoring.
Overall, the non-parametric and adaptive characteristics of KDE provide crucial support for lightweight modeling and robustness enhancement in consumer electronic systems; however, its applications in high-dimensional and complex scenarios remain limited and warrant further investigation.


To improve the computational efficiency of high-dimensional KDE, researchers have mainly focused on three optimization directions: space partition, dimensionality reduction, and hash-based retrieval.
Space partition divides the dataset $P$ into several subsets $P_l$, using local $\mathrm{KDE}_{P_l}(q)$ to replace global computation and reduce complexity. \cite{reff6} employed KD-tree hierarchical sampling combined with Monte Carlo approximation to enhance computational efficiency, though it requires a large number of samples under high-variance distributions. \cite{reff8} proposed a tree-structured partition method for road networks to achieve dynamic spatiotemporal KDE, where the partition granularity must balance between accuracy and efficiency.
In terms of dimensionality reduction, \cite{reff9} proposed DC-KDE, which integrates deep autoencoders with geometric correction for nonlinear dimensionality reduction, effectively improving density estimation accuracy for manifold data; \cite{reff10} applied PCA for dimensionality reduction followed by KDE and used inverse mapping to correct bias, though dimensionality reduction may distort the original structure.
Hash-based retrieval provides another efficient approach. Traditional methods \cite{reff13,reff15,reff17} such as LSH or PQ structures approximate kernel value computation through approximate nearest-neighbor search, significantly accelerating KDE; however, LSH and PQ fails to capture hierarchical semantic features. Recently, deep hash methods such as UHSCM\cite{tu2023unsupervised}, SDC\cite{ng2023unsupervised}, FSCH\cite{cao2023unsupervised}, HARR\cite{ma2024harr}, and UATP\cite{tang2024once} have improved partition quality and retrieval efficiency by learning discriminative hash codes，mainly used for images and videos. However, while these methods are mainly designed for images and videos, they cannot be directly adapted to high-dimensional vectors. Inspired by these advances, we introduce deep hash strategies into parameter optimization and spatial partition to achieve more efficient high-dimensional KDE optimization and modeling.
\section{Multi-Encoder Contrastive Hashing}

\subsection{Network Structure}
为了克服传统自编码哈希的局限性，受变分自编码器（VAE）和对比学习的启发，我们提出了一种现代无监督哈希学习框架，称之为多编码器对比哈希（Multi-Encoder Contrastive Hashing，简称 MECH）。与采用单一编码器-解码器架构的传统基于 VAE 的哈希方法不同，本框架采用了“多编码器多哈希表”的结构。其中，多个独立的编码器用于为不同的哈希表生成多样化的哈希表示，而特征重构则由一个共享的解码器共同完成。这种设计使模型能够学习到互补的哈希投影，并在多个哈希表上执行并行检索，从而显著提升了检索精度与鲁棒性。同时，通过在不同编码器之间共享解码器，该架构大幅减少了参数量与重构开销，从而提高了训练效率并降低了存储复杂度
To overcome the limitations of traditional autoencoding hashing, we propose a modern unsupervised hashing learning framework inspired by variational autoencoders (VAE) and contrastive learning,termed Multi-Encoder Contrastive Hashing(MECH). Different from conventional VAE-based hashing methods that employ a single encoder-decoder architecture, the proposed framework adopts a multi-encoder multi-hash-table structure, where multiple independent encoders are used to generate diverse hash representations for different hash tables, while a shared decoder is employed for feature reconstruction. This design enables the model to learn complementary hash projections and perform parallel retrieval across multiple hash tables, thereby improving retrieval accuracy and robustness. Meanwhile, by sharing the decoder among different encoders, the proposed architecture significantly reduces the number of parameters and reconstruction overhead, leading to improved training efficiency and lower storage complexity.

The architecture consists of three main components:

1) Multi-Encoder: maps input vector $\mathbf{x}$ into a low-dimensional probabilistic latent space, producing mean $\boldsymbol{\mu}$ and variance $\boldsymbol{\sigma}^2$ for the latent variable $\mathbf{z}$:
\[
\mathbf{z} \sim \mathcal{N}(\boldsymbol{\mu}, \text{diag}(\boldsymbol{\sigma}^2))
\]

2) Hash Layer: applies a continuous relaxation to generate approximately binary hash codes:
\[
\mathbf{h} = \tanh(\mathbf{z})
\]

3) Shared-Decoder: reconstructs input from latent variable $\mathbf{z}$ to ensure information preservation:
\[
\hat{\mathbf{x}} = \text{Shared-Decoder}(\mathbf{z})
\]

4) Multi-table integration with a shared decoder: This approach trains multiple independent encoder hash models while employing a single shared decoder. By feeding multiple hash projections into the shared decoder, it generates several hash tables simultaneously. Compared with constructing each hash table using fully independent models, this design substantially reduces both training complexity and storage requirements.
\begin{equation*}
k_\sigma(x,y) = \exp\!\left( -\frac{\|x-y\|^2}{\sigma^2} \right)
\end{equation*}
\begin{figure*}
    \centering
    \includegraphics[width=1\linewidth]{image/架构图2.png}
    \caption{Structure diagram of Hashing with Contrastive Learning}
    \label{fig:mech_architecture}
\end{figure*}



\subsection{Loss Functions}

The proposed framework jointly optimizes the following losses:

Reconstruction Loss $L_{rec}$: 
Ensures that latent representations preserve input information:
\[
L_{rec} = \mathbb{E}_{q(\mathbf{z}|\mathbf{x})} \big[ \|\mathbf{x} - \hat{\mathbf{x}}\|^2 \big]
\]

Contrastive Similarity Loss $L_{contrast}$: To preserve the neighborhood structure of the original feature space, we formulate a contrastive objective based on pairwise relative distance relationships.

Given an anchor sample $\mathbf{x}_i$, we construct a positive pair $\mathbf{x}_j$ and a negative sample $\mathbf{x}_k$ according to their Euclidean distances, where:
\[
\|\mathbf{x}_i - \mathbf{x}_j\|_2 < \|\mathbf{x}_i - \mathbf{x}_k\|_2.
\]

The corresponding hash codes are expected to preserve the same similarity ordering in the Hamming space:
\[
D_H(\mathbf{b}_i, \mathbf{b}_j) < D_H(\mathbf{b}_i, \mathbf{b}_k),
\]
where $D_H(\cdot,\cdot)$ denotes the Hamming distance between binary hash codes.

To enforce this constraint, we define the following margin-based contrastive similarity loss:
\[
L_{contrast} =
\sum_{(i,j,k) \in \mathcal{T}}
\max\Big(0, m + D_H(\mathbf{b}_i, \mathbf{b}_j) - D_H(\mathbf{b}_i, \mathbf{b}_k)\Big),
\]
where $m$ is a margin hyper-parameter, and $\mathcal{T}$ denotes the set of sampled triplets constructed from Euclidean distance ranking.

The Hamming distance is defined as:
\[
D_H(\mathbf{b}_i, \mathbf{b}_j)
=
\frac{1}{2}\left(c - \mathbf{b}_i^\top \mathbf{b}_j\right),
\]
where $c$ is the hash code length.

This loss encourages the learned hash space to preserve the relative similarity structure of the original feature space, ensuring that closer samples in Euclidean space remain closer in Hamming space, while distant samples are pushed further apart.

Smooth Quantization Loss $L_{quan}$: 
Constrains hash codes to approach binary values while maintaining gradient flow:
\[
\mathbf{b}=\operatorname{sign}(\mathbf{h}), \qquad
L_{quan} = \|\mathbf{h} - \mathbf{b}\|^2
\]

To avoid degenerated hash tables, we additionally use bit-balance and bit-decorrelation regularization:
\[
L_{bal}=\left\|\frac{1}{n}\sum_{i=1}^{n}\mathbf{h}_i\right\|_2^2,\qquad
L_{dec}=\left\|\frac{1}{n}\mathbf{H}^{\top}\mathbf{H}
-\operatorname{diag}\left(\frac{1}{n}\mathbf{H}^{\top}\mathbf{H}\right)\right\|_F^2 .
\]

Total Loss:
The final objective balances reconstruction, latent regularization, contrastive similarity, and quantization:

\begin{equation}
L =  L_{rec} + \alpha L_{contrast} + \beta L_{quan}
    + \gamma L_{bal} + \eta L_{dec}
\label{eqL}
\end{equation}
where the weighting coefficients $\alpha$, $\beta$, $\gamma$, and $\eta$ control the relative contribution of contrastive similarity, quantization, bit balance, and bit decorrelation, respectively. By default, we set $\alpha = 1$, $\beta = 0.1$, $\gamma = 0.1$, and $\eta = 0.1$ to balance semantic similarity preservation, binary regularization, and hash-table utilization. Experimental results demonstrate that the proposed framework is not sensitive to the choice of these coefficients within a broad range of values in Table~\ref{tab:sensitivity_cifar10_small}.



\subsection{Training and Retrieval}
为确保计算效率，本研究中变分自编码器（VAE）的基础映射单元采用轻量化的对称多层感知机（MLP）架构，编码与解码阶段均仅包含单层全连接映射。其中，编码器的输入维度取决于具体数据集的高维特征 $\mathbf{x}$，随后通过两个并行的全连接层分别输出均值 $\boldsymbol{\mu}$ 与方差 $\boldsymbol{\sigma}^2$，其维度与潜变量 $\mathbf{z}$ 的维数一致。该潜空间进一步通过哈希层 $\tanh(\alpha \mathbf{z})$ 进行连续松弛化映射以控制输出的码长。
在实验流程与超参数设置方面，各数据集均被严格划分为独立的训练集与测试集，且该划分在后续对比实验中保持一致，以杜绝数据泄露风险。为了在计算效率与检索精度之间取得平衡，哈希码长度 $c$ 分别设置为 16、32、64 和 128 位。整个网络采用 Adam 优化器进行端到端训练，基础学习率设为 0.001，批处理大小（Batch Size）为 128，总迭代轮数（Epochs）为 100。


To ensure computational efficiency, the proposed MECH framework adopts a lightweight symmetric multi-layer perceptron (MLP) architecture as the basic mapping unit for both encoding and decoding stages. Specifically, each encoder and the shared decoder contain only a single fully connected layer.

For an input high-dimensional feature vector $\mathbf{x}$, the encoder first maps the input into a latent probabilistic space through two parallel fully connected layers that separately estimate the mean vector $\boldsymbol{\mu}$ and variance vector $\boldsymbol{\sigma}^2$. The dimensionality of both vectors is identical to that of the latent variable $\mathbf{z}$. The latent representation is then transformed into continuous hash codes through the hash layer.
\[
\mathbf{h} = \tanh(\mathbf{z})
\]


To improve representation diversity and retrieval robustness, the proposed framework employs multiple independent encoders to construct multiple complementary hash tables, while a shared decoder is utilized to reconstruct input features from different latent representations. Compared with fully independent encoder-decoder architectures for each hash table, the shared-decoder design significantly reduces parameter redundancy and reconstruction overhead.

During training, each dataset is strictly divided into independent training and testing subsets, and the same data partition is consistently adopted across all comparison experiments to avoid data leakage. To balance retrieval accuracy and computational efficiency, the hash code length $c$ is set to 16, 32, 64, and 128 bits, respectively.

The entire framework is optimized end-to-end using the Adam optimizer with an initial learning rate of 0.001. The batch size is set to 128, and the total number of training epochs is 100. For contrastive learning, two augmented views of the same sample are treated as positive pairs, while all remaining samples within the mini-batch are regarded as negative pairs. The temperature parameter in the contrastive loss is empirically set to $\tau = 0.5$.

During retrieval, query samples are simultaneously processed by all encoder branches to generate multiple binary hash codes.Each binary code independently indexes its corresponding hash table, and the retrieved candidate sets from different hash tables are subsequently merged for final approximate nearest neighbor search.

The Hamming distance between binary hash codes is computed as:
\[
D_H(\mathbf{b}_i,\mathbf{b}_j)
=
\frac{1}{2}
\left(
c-\mathbf{b}_i^\top \mathbf{b}_j
\right)
\]
where $c$ denotes the hash code length.

During training, random data augmentations are applied to create positive pairs for contrastive learning. The encoder produces latent codes $\mathbf{z}$, which are mapped to hash codes $\mathbf{h}$. Decoder reconstructs inputs to maintain information integrity. For retrieval, the query is encoded into hash codes $\mathbf{h}_q$, which are then matched against hash codes in the table(s) using Hamming distance or inner-product similarity. This framework removes the need for angle-based similarity calculations, ensures smooth binarization, and leverages modern self-supervised learning to improve hash code quality and retrieval performance in high-dimensional spaces.


\section{Kernel Density Estimation Based on Approximate Annulus Partition}
In the following section, we first introduce the method based on accurate annulus partition and, on this basis, propose a more practical kernel density estimation approach using approximate annulus partition.

\subsection{Kernel Density Estimation Based on Accurate Annulus partition}
In this section, we define the concept of accurate annulus partition, which lays the theoretical foundation for the design of the approximate annulus partition method.
 

Let $M = (X, \mathrm{dis})$ denote a metric space.
Given $p, q \in X$, the Euclidean distance is defined as ${\mathrm{dis}} = |p - q|_2$.
Unless otherwise specified, the $l_2$ norm is expressed as $|\cdot|$.
We define the closed ball centered at point $q$ with radius $r$ as:
\begin{equation*}
B(q,r) = \{ p \in P \mid \mathrm{dis}(q,p) \leq r \}
\end{equation*}


For a point set $P \subset \mathbb{R}^d$ containing $n$ points, and a kernel function $k(\cdot,\cdot): \mathbb{R}^d \times \mathbb{R}^d \to [0,1]$,
the kernel density function of any point $x \in \mathbb{R}^d$ with respect to the dataset $P \subset \mathbb{R}^d$ is defined as:

\begin{equation*}
\mathrm{KDE}_P(x) = \frac{1}{n} \sum_{y \in P} k(x,y)
\end{equation*}


Given a query point $q$, the kernel density can be estimated using a sampling approach as $\mathrm{KDE}_P(q) = \frac{1}{|P|} \sum_{p \in P} k(q,p)$.
If the dataset $P$ is not specifically indicated, we simply denote it as $\mathrm{KDE}(q)$.
Let $\Phi(q,P) = \sum_{p \in P} k(q,p)$, that is, $\mathrm{KDE}(q) = \frac{\Phi(q,P)}{n},$
where $n$ represents the number of elements in the dataset $P$.This paper mainly focuses on the widely used Gaussian kernel density.

\begin{equation*}
k_\sigma(x,y) = \exp\!\left( -\frac{\|x-y\|^2}{\sigma^2} \right)
\end{equation*}
\begin{figure*}
    \centering
    \includegraphics[width=1\linewidth]{image/架构图2.png}
    \caption{Structure diagram of kernel density estimation based on accurate annulus partition}
    \label{fig:accurate_annulus}
\end{figure*}

An annulus $A$ can be determined when a query point $q$ and radii $r_1, r_2$ are given.
$A$ refers to all points in $P \subset \mathbb{R}^d$ whose distances from $q$ lie within the interval $[r_1, r_2]$, denoted as $A = A_q(r_1, r_2)$.
When $|r_1 - r_2|$ is small, the distances $d(q, p)$ of all points $p \in A_q$ differ only slightly and are of the same order of magnitude, resulting in minimal variation of corresponding kernel values.
In this case, the kernel density at the query point can be approximately estimated by counting the number of points in each annulus, while assigning sampling weights to different annuli to reduce the number of sampling points.
Compared with other partition methods, the annular structure can cover the data space with fewer partitioned subsets while maintaining small kernel density variation within each annulus, thereby significantly improving computational efficiency under controlled error conditions.

The dataset $P$ is partitioned into annuli $A_1, A_2, \ldots, A_\phi$ centered at $q$.
Let $n_i = |A_i|$ denote the number of elements in partition $A_i$, with $\sum_{i=1}^{\phi} n_i = n$.
Let $S_i$ be an independent random sampling set of partition $A_i$, and the number of samples be $m_i = |S_i|$.
We define the kernel density random variable of a sampled point $x$ as $X_j = k(q,x)$, and define the kernel density estimation of the accurate annulus partition $A_i$ as:

\begin{equation*}
Z_{A_i} = \frac{n_i}{m_i} \sum_{j=1}^{m_i} X_j
\end{equation*}

We obtain the kernel density estimate for the accurate annulus partition centered at the query point $q$:

\begin{equation*}
\Phi(q,P) = \sum_{p \in P} k(q,p) = \sum_{i=1}^{\phi}Z_{A_i}
\end{equation*}

Efficient annulus-based kernel density estimation relies on two key components: 1) an annulus partition strategy, which must ensure both structural rationality and query dependency to efficiently form hierarchical annuli centered on the query point; 2) an annulus density estimation design, which determines the number of samples and serves as the core factor influencing estimation efficiency.

However, the accurate annulus partition inevitably requires computing the distance between the query point and all data points, which is equivalent to a full kernel density computation and inefficient.
To address this, we propose in the following section an 
approximate annulus partition scheme. Within the same annulus, a small number of samples can represent the overall distribution, achieving a balance between the number and quality of subsets while maintaining both estimation accuracy and computational efficiency.

\subsection{Angular Space Approximate Nearest Neighbor Query}
\begin{figure*}
    \centering
    \includegraphics[width=1\linewidth]{image/近似环划分原理图.png}
    \caption{Schematic diagram of the approximate annulus partition}
    \label{fig:approx_annulus}
\end{figure*}

During approximate annulus partition, the MECH method is first used to transform the annulus into an angular space approximate nearest neighbor query.

\textbf{Hash Tables}:
Given a family of hash functions $G$, to construct $L$ hash tables, we randomly and independently select $L$ functions $g_1, g_2, \ldots, g_L$ from $G$.
For each $l = 1, 2, \ldots, L$, the point $p \in P$ is stored in the hash bucket $HT_l$, and all hash buckets together form the hash table structure. 

In a space where the angle is used as the distance metric, the MECH method ensures that points with smaller angles have a higher probability of falling into the same hash bucket, or that the Hamming distance between the corresponding hash values of two points, $|g(q) - g(x)|$, is small.


Given a query point $q$, we denote the accumulated collision probability as:
\begin{equation*}
P(x) = \sum_{i \in{J}} P(d_{qx} = i \mid \theta_{qx})
\end{equation*}
where $\theta_{qx}$ represents the angle between $x$ and $q$, $\theta_{qx} \in [0,\pi]$, and ${J}$ denotes the distance set, ${J} = {0,1,\ldots,I}$.
Thus, in a single hash mapping $g(v)$, the probability that a point with an angle $\theta_{qx}$ from $q$ has a Hamming distance $d_{qx} \in {J}$ between its hash value and that of $q$ is $p(x)$.
Accordingly, in $L$ hash tables, the probability that a point with an angle $\theta_{qx}$ from $q$ has its hash value satisfying $d_{qx} \in {J}$ in at least one hash table is:
\[
1 - (1 - p)^L
\]  
We expect this probability to be greater than $\alpha$:
\[
1 - (1 - p)^L > \alpha
\]  
To determine $I$, the minimum value of $I$ must satisfy:

\begin{equation}
\begin{aligned}
I = \min_{I \in [0,k]} \Bigg(
    &\sum_{i=0}^I P(d_{qx} = i \mid \theta_{qx}) \\
    &> 1 - (1 - \alpha)^{\tfrac{1}{L}}
\Bigg)
\end{aligned}
\label{eqI_threshold}
\end{equation}



\textbf{Distance Tables}: 
For each hash value $g^l$ in the hash tables, where $l \in {1,2,\ldots,L}$, we define $DT_h^l(i)$ as the set of all hash values in the $l$-th table whose hash value is $h$ and whose hamming distance is $i$, where $i = 0, 1, \ldots, k$.
The key to efficiently retrieving relevant points during the query phase lies in the distance tables.

Given a query angle $\theta$, according to \eqref{eqI_threshold}, we can determine the distance set $J$.
From the distance tables, all hash values within the distance range corresponding to the query point are retrieved.
By combining these with the hash tables, the points corresponding to the retrieved hash values are collected as the query candidate set, denoted as $C_q(\theta)$:

\begin{equation*}
C_q(\theta) = \bigcup_{l=1}^L \bigcup_{i \in J} HT_l \big( DT_h^l(i) \big)
\end{equation*}

\subsection{Hierarchical Distance Transformation}
We have already demonstrated how to perform approximate nearest neighbor queries in angular space using MECH. It is further necessary to explain how to transform annuli from euclidean space to angular space, where the main concept applied is hierarchical distance transformation.


Given two points $p, q \in \mathbb{R}^d$,
to achieve the transformation from euclidean space to angular space, we first apply the cosine theorem, under which the angular relationship between $p$ and $q$ with respect to the origin satisfies:
\begin{equation*}
 d_{pq} = \|p - q\|
\end{equation*}


\begin{equation*}
\cos(\theta_{pq}) = \frac{\|q\|^2 + \|p\|^2 - d_{pq}^2}{2 \times \|q\| \times \|p\|} 
\end{equation*}

According to the above formula, when $q$ and $r$ are known, the angular range of points whose distances from $q$ fall within the interval $[0, r]$ can be expressed as：

\begin{equation*}
\begin{aligned}
\theta_{B_q(r)} = \arccos \left( \frac{\|q\|^2 + \|p\|^2 - r^2}{2 \times \|q\| \times \|p\|} \right) 
\end{aligned}
\end{equation*}


At this point, the maximum angle between the points in $B_q(r)$ and the query point $q$ is:


\begin{equation}
\theta_{\max} = \arcsin \left( \frac{r}{\|q\|} \right)
\label{eq01}
\end{equation}

When using MECH to query points whose angles with $q$ are approximately within $\theta_{\max}$, the results may exhibit significant errors.
This occurs particularly when the norm of a point $|p|$ does not satisfy the conditions, such as in regions where $|p| < |q| - r$ or $|p| > |q| + r$.
To reduce this error, we can partition the dataset into spherical layers to form a sphere table.




\textbf{Sphere Table}:
Given a dataset $P$ and $\delta \in \mathbb{R}_+$, we define the layer index for each $p \in P$ as $s = \mathrm{round}\left(\frac{|p|}{\delta}\right)$ and store it in the Sphere Table indexed by $s$, denoted as $ST_s$. Through this hierarchical processing, the dataset is partitioned into concentric layers centered at the origin, with an interval of $\delta$.

When a query point $q$ and query radius $r$ are given, we can determine the layer range involved in the approximate nearest-neighbor query $B_q(r)$. The minimum layer index corresponding to $B_q(r)$ is $fs = \mathrm{round}\!\left(\frac{|q| - r}{\delta}\right),$ and the maximum layer index is $ls = \mathrm{round}\!\left(\frac{|q| + r}{\delta}\right)$, and we denote the set of points within the query radius $r$ around point $q$ as $ST_{q,r}$, $ST_{q,r}= \bigcup_{fs}^{ls} ST_s$. At this point, compared with querying points whose angles with $q$ are approximately within $\theta_{\max}$, we can compute for each annulus the maximum angle $\theta_{\max}$ between $[fs, ls]$. we can denote the approximate nearest-neighbor query $B_q(r)$ result as:
\[
B_q(r) = ST_{q,r} \cap C_q(\theta_{\max})
\]


\subsection{Kernel Density Estimation Based on Approximate Annulus Partition}
The main idea of the approximate annulus query method is to utilize MECH to construct two filtering approaches consisting of angular query and hierarchical distance transformation for performing approximate annulus partition. Unlike conventional hash methods in Euclidean space, this method ensures that the annulus parameters $r_1$ and $r_2$ can be specified during the query stage.


In the preprocessing stage, since the algorithm is primarily query-driven, only a single preprocessing operation is required to handle different query tasks. During this stage, three data structures—Sphere Table, Hash Tables, and Distance Tables—are constructed sequentially and stored using dictionary structures. In the query stage, indexing can be performed directly without additional computation.


In the query stage, given a query point $q$ and the inner and outer annulus radii $r_1$ and $r_2$, for $r_1$, we can compute the maximum angle $\theta_{\max}$ of the points in each layer using \eqref{eq01}. According to \eqref{eqI_threshold}, the corresponding distance set $J$ can then be obtained. Finally, by retrieving all hash values from the Distance Tables whose hamming distances from the hash value of $q$ within the query radius $r_1$ belong to $J$, and combining them with the Sphere Table for distance transformation, we obtain the final approximate query point set $B_q(r_1)$.
Similarly, for a given $r_2$, we can obtain the approximate query result $B_q(r_2)$. Since $r_2 < r_1$, it follows that $B_q(r_2) \subset B_q(r_1)$. Following the standard procedure, after obtaining $B_q(r_1)$ and $B_q(r_2)$, we can simply compute their difference to obtain the approximate annulus query result, as shown in Fig.~\ref{fig:approx_annulus}, denoted as: 
\[
A_q(r_1,r_2) = B_q(r_1) \setminus B_q(r_2)
\]  
The advantage of this method lies in its use of a difference set strategy, which effectively removes points that are highly likely to belong to $B_q(r_2)$, while ensuring that points in $A_q(r_1,r_2)$ are not erroneously excluded due to a low-probability collision with $B_q(r_2)$ in any single hash table. This approach improves recall while maintaining a reasonable level of accuracy.Although distance transformation adopts the idea of trading a small amount of accuracy for a feasible solution, the transformed angular space offers unique advantages. The angular space can also compensate for certain limitations inherent in Euclidean space. For instance, MECH used for approximate nearest neighbor search in angular space benefits from the bounded nature of its hash values ($2^k$), whereas euclidean-space hash codes are unbounded, making them more difficult to process and store. This characteristic was the initial motivation for applying this concept to address approximate annulus partition in euclidean space since unbounded hash codes in euclidean space make multi-probe searches over wide ranges infeasible. By leveraging the advantages of angular space, the distance transformation concept provides an effective solution to the approximate annulus partition problem.


The approximate annulus-based kernel density estimation process is similar to that of the accurate annulus-based approach. Centered on the query point $q$, the dataset is divided into several approximate annulus $A_1, A_2, \ldots, A_\phi$. Since the distances between points within each annulus and $q$ are similar and the kernel function varies only slightly, each annulus can be regarded as a homogeneous region. By performing independent random sampling within each annulus and computing a weighted sum based on the number of points in each ring, the kernel density at the query point can be efficiently approximated. Compared with traditional uniform or grid-based partition methods, annulus partition significantly reduces computational complexity and sampling cost while maintaining estimation accuracy.
\subsection{Influence of Recall and Precision on KDE Accuracy}

The quality of approximate annulus retrieval directly affects the accuracy of kernel density estimation. 
For approximate annulus partition, the retrieved annulus $\hat{A}_i$ may contain two types of retrieval errors: 
1) false negatives, where relevant points belonging to the true annulus are not retrieved; 
2) false positives, where irrelevant points outside the target annulus are mistakenly included. 
These two errors correspond to recall deficiency and precision deficiency, respectively.

Let $A_i$ denote the exact annulus and $\hat{A}_i$ denote the approximate annulus obtained by MECH retrieval. 
The recall and precision are defined as:

\[
\mathrm{Recall}
=
\frac{|A_i \cap \hat{A}_i|}{|A_i|}
\]

\[
\mathrm{Precision}
=
\frac{|A_i \cap \hat{A}_i|}{|\hat{A}_i|}
\]

For Gaussian kernel density estimation, the kernel function is:

\[
k(q,p)
=
\exp\left(
-\frac{\|q-p\|^2}{\sigma^2}
\right)
\]

where the kernel contribution decreases exponentially as the distance increases. 
Therefore, points close to the query point contribute significantly larger kernel values than distant points.

The exact kernel density contribution of annulus $A_i$ is:

\[
\Phi_i(q)
=
\sum_{p \in A_i} k(q,p)
\]

while the approximate annulus estimation is:

\[
\hat{\Phi}_i(q)
=
\sum_{p \in \hat{A}_i} k(q,p)
\]

The estimation error can then be expressed as:

\[
\hat{\Phi}_i(q)-\Phi_i(q)
=
\sum_{p\in \hat{A}_i\setminus A_i}k(q,p)
-
\sum_{p\in A_i\setminus \hat{A}_i}k(q,p)
\]

The first term corresponds to the error introduced by false-positive points, while the second term corresponds to the error caused by false-negative points.

For false-positive points, their distances to the query point are usually larger than the target annulus radius. 
Due to the exponential decay property of the Gaussian kernel, these points generally satisfy:

\[
k(q,p)\approx 0
\]

Thus, although low precision introduces additional irrelevant points, their kernel contributions are typically very small and only slightly affect the final density estimation.

In contrast, false-negative points are true neighboring points whose distances satisfy:

\[
\|q-p\|\in[r_1,r_2]
\]

These points usually possess relatively large kernel contributions. 
Missing such points directly reduces the estimated density value and introduces significant estimation bias.

Assume the average kernel contribution of true annulus points is $\bar{k}_{in}$, and the average contribution of false-positive points is $\bar{k}_{out}$. 
In general:

\[
\bar{k}_{in}\gg \bar{k}_{out}
\]

Then the expected KDE error approximately satisfies:

\[
|\hat{\Phi}_i(q)-\Phi_i(q)|
\approx
N_{FN}\bar{k}_{in}
+
N_{FP}\bar{k}_{out}
\]

where $N_{FN}$ and $N_{FP}$ denote the numbers of false-negative and false-positive points, respectively.

Since $\bar{k}_{in}\gg \bar{k}_{out}$, the influence of false negatives dominates the overall KDE estimation error. 
Therefore, recall plays a more important role than precision in annulus-based kernel density estimation. 
A higher recall ensures that most true neighboring points are preserved, thereby maintaining the stability and accuracy of density estimation, whereas moderate precision reduction usually introduces only limited additional error.

This analysis explains why the proposed MECH-based approximate annulus partition can reduce KDE estimation error when high-contribution points are preserved. 
Although the point-level precision of MECH is not always the highest across all datasets, the kernel-weighted metrics in the experiments show whether the retrieved annulus preserves the contributions that dominate KDE estimation.

\section{Experiments}
This study evaluates the MECH-based approximate annulus KDE framework through annulus retrieval, KDE estimation, structural ablation, parameter sensitivity, and downstream anomaly detection.

\subsection{Experimental Design}
We conduct annulus partition experiments on GIST-512, CIFAR-10-Small, Amazon-Commerce-Reviews, and ISOLET. These datasets differ in dimension, sample size, and feature type, and are used to test whether approximate annulus construction remains stable under heterogeneous high-dimensional settings. For the additional MECH error-propagation experiment reported in Tables~\ref{tab:mech_hash_comparison}--\ref{tab:mech_sensitivity}, we use ISOLET with 2800 training points, 1800 index points, and 40 query points. The base annulus is defined by the 15th--25th percentiles of the query-to-index distance distribution, and the base Gaussian bandwidth is 6.7226. MECH is trained with 12 encoder branches and a 12-bit maximum code length for 24 epochs; unless otherwise specified, the annulus index uses $L=12$ hash tables, $c=4$ hash bits, $\delta=0.25$, and Hamming probe radius 1.

For kernel density anomaly detection, we use UNSW-NB15, which contains 49 network traffic features and a binary label column. Each dataset is split into independent training, validation, and test partitions. The MECH model is trained only on the training partition. Validation data are used to select the Hamming threshold, loss weights, kernel bandwidth, and anomaly detection threshold. Reported metrics are computed on the test partition. Before constructing the sphere table, each dataset is first normalized by min-max scaling and then moved farther away from the origin. Specifically, we select the normalized sample with the largest Euclidean norm as an anchor vector $\mathbf{a}$ and transform every vector as $\mathbf{x}'=2\mathbf{x}+\mathbf{a}$. This transformation uniformly scales pairwise Euclidean distances by a factor of two while preserving their relative ordering, and it prevents the data distribution from concentrating near the origin, which makes the origin-centered sphere layers and angular distance transformation more stable.

\subsection{Kernel-Weighted Evaluation Metrics}
Point-level precision, recall, and F1-score indicate whether the retrieved annulus contains the correct points, but KDE depends on kernel contributions rather than point counts. Therefore, for an exact annulus $A_i$, an approximate annulus $\hat{A}_i$, query point $q$, and Gaussian kernel value $k(q,x)$, we additionally report:
\[
R_i^k = \frac{\sum_{x\in A_i\cap \hat{A}_i} k(q,x)}
{\sum_{x\in A_i} k(q,x)+\epsilon},
\]
\[
E_{+,i}^k = \frac{\sum_{x\in \hat{A}_i\setminus A_i} k(q,x)}
{\sum_{x\in A_i} k(q,x)+\epsilon},\qquad
E_{-,i}^k = \frac{\sum_{x\in A_i\setminus \hat{A}_i} k(q,x)}
{\sum_{x\in A_i} k(q,x)+\epsilon}.
\]
The annulus KDE relative error is computed as
\[
e_{\mathrm{KDE}} =
\frac{\left|\sum_{x\in \hat{A}_i} k(q,x)-\sum_{x\in A_i} k(q,x)\right|}
{\sum_{x\in A_i} k(q,x)+\epsilon}.
\]
We also report the candidate expansion ratio (CER), defined as $|\hat{A}_i|/(|A_i|+\epsilon)$, and the online query time. The reported time excludes MECH training and the preprocessing construction of Hash Tables, Distance Tables, and Sphere Tables. In this setting, MECH query hash codes are cached and stored after preprocessing, whereas SimHash, Angular LSH, and Hyperplane LSH compute query hash codes online because they do not use cached query hash tables. The reported online time therefore measures the actual KDE query path under this deployment setting, including candidate retrieval, candidate filtering, and Gaussian kernel accumulation over the final approximate annulus candidates. For timing stability, each query is measured seven times and the median runtime is used before averaging over all queries.

\subsection{MECH Loss-Weight Sensitivity}
\begin{table}[htbp]
\centering
\caption{Sensitivity analysis of the weighting coefficients on CIFAR-10-Small (\% F1-score).}
\label{tab:sensitivity_cifar10_small}
\begin{tabular}{c|ccccc}
\hline
$\beta$ & $\alpha=0.1$ & $\alpha=0.5$ & $\alpha=1$ & $\alpha=5$ & $\alpha=10$ \\
\hline
0.001 & 76.4 & 78.1 & 79.0 & 78.3 & 77.1 \\
0.01  & 77.3 & 79.2 & 80.1 & 79.5 & 78.0 \\
0.1   & 78.0 & 80.0 & \textbf{80.8} & 80.2 & 79.1 \\
1     & 75.8 & 77.0 & 78.2 & 77.4 & 76.3 \\
\hline
\end{tabular}
\end{table}

Table~\ref{tab:sensitivity_cifar10_small} compares different settings of the weighting coefficients $\alpha$ and $\beta$ in Equation~\eqref{eqL}. The best performance is achieved when $\alpha=1$ and $\beta=0.1$, which is used in the remaining experiments together with $\gamma=\eta=0.1$ for bit balance and decorrelation.

\subsection{Approximate Annulus Partition}
This section evaluates approximate annulus partition using SimHash, Angular LSH, Hyperplane LSH, and MECH. All methods are integrated with the same sphere-layer radial filtering and difference-based annulus construction. Figures~\ref{fig:precision}--\ref{fig:F1} show the original point-level comparison under varying numbers of hash tables.

\begin{figure}
    \centering
    \includegraphics[width=1\linewidth]{image/precision.png}
    \caption{Precision of approximate annulus partition with varying numbers of hash tables on multiple datasets.}
    \label{fig:precision}
\end{figure}

\begin{figure}
    \centering
    \includegraphics[width=1\linewidth]{image/recall.png}
    \caption{Recall of approximate annulus partition with varying numbers of hash tables on multiple datasets.}
    \label{fig:recall}
\end{figure}

\begin{figure}
    \centering
    \includegraphics[width=1\linewidth]{image/F1.png}
    \caption{F1-score of approximate annulus partition with varying numbers of hash tables on multiple datasets.}
    \label{fig:F1}
\end{figure}

The point-level curves show that learned hashing is competitive, but point-level F1 alone cannot determine KDE accuracy. Table~\ref{tab:mech_hash_comparison} therefore reports kernel-weighted annulus metrics on ISOLET after the origin-shift preprocessing. To avoid hash-candidate saturation, the hash backbone comparison uses a non-saturated query setting with 4 hash tables, 12-bit codes, and exact-bucket lookup. Under this setting, random-projection hash backbones miss a larger fraction of true annulus points, while MECH preserves substantially more kernel-contributing candidates and yields the lowest KDE relative error.

\begin{table}[htbp]
\centering
\caption{Hash method comparison with kernel-weighted annulus KDE metrics on ISOLET.}
\label{tab:mech_hash_comparison}
\resizebox{\linewidth}{!}{%
\begin{tabular}{lcccccccc}
\toprule
Method & P & R & F1 & $R^k$ & $E_+^k$ & KDE Err. & CER & Online ms \\
\midrule
SimHash & 96.86\% & 89.01\% & 92.09\% & 89.03\% & 2.89\% & 8.49\% & 0.92 & 0.680 \\
Angular LSH & 96.93\% & 89.99\% & 93.18\% & 89.99\% & 2.86\% & 7.15\% & 0.93 & 0.675 \\
Hyperplane LSH & 96.95\% & 87.32\% & 91.65\% & 87.34\% & 2.74\% & 9.96\% & 0.90 & 0.679 \\
MECH & 96.89\% & 96.75\% & 96.82\% & 96.76\% & 3.10\% & 1.43\% & 1.00 & 0.648 \\
\bottomrule
\end{tabular}%
}
\end{table}

\subsection{KDE Estimation}
This section compares accurate annulus partition, approximate annulus partition with MECH, and Monte Carlo Markov Chain (MCMC) sampling. The evaluation considers both the number of sampled points required under fixed relative error and the relative error under fixed sample size.

Fig.~\ref{fig:error} shows the average number of sampled points when relative errors are fixed at 0.01, 0.03, and 0.05. Accurate annulus partition requires the fewest samples, while approximate annulus partition achieves comparable performance with lower indexing cost. Both annulus-based methods require substantially fewer samples than MCMC.

\begin{figure}
    \centering
    \includegraphics[width=1\linewidth]{image/relative error1.png}
    \caption{Sampling counts among accurate annulus, approximate annulus, and MCMC under fixed relative error.}
    \label{fig:error}
\end{figure}

Fig.~\ref{fig:Sampling} compares relative errors under different sampling sizes. The approximate annulus method remains close to accurate annulus partition, especially on CIFAR-10-Small and ISOLET, and both annulus methods stabilize after the sample size exceeds 200.

\begin{figure}
    \centering
    \includegraphics[width=1\linewidth]{image/number of sample1.png}
    \caption{Relative error comparison among accurate annulus, approximate annulus, and MCMC across datasets.}
    \label{fig:Sampling}
\end{figure}

\subsection{Structure Ablation}
Table~\ref{tab:mech_structure_ablation} evaluates the contribution of each deployable retrieval component. In this ablation, all variants use the same cached MECH hash candidates, so the reported online time includes materializing the cached candidate list, applying the variant-specific filtering operations, and computing the Gaussian KDE sum over the final candidates. Hash-only retrieval has little filtering cost, but it must evaluate the kernel over many false-positive candidates and therefore introduces excessive weighted false positives and candidate expansion. Adding only sphere filtering cannot define a valid query-centered annulus because radial layers are centered at the origin and still require angular constraints. The distance-table-only variant uses a global angular threshold without sphere-layer radial constraints, so it misses many true annulus points and produces larger KDE error. Full MECH combines hash retrieval, distance-table angular filtering, sphere-layer radial filtering, and difference-based annulus construction. It therefore adds filtering cost, but reduces the final KDE candidate set and achieves the best practical structure result: the highest precision, recall, kernel-weighted recall, and the lowest KDE relative error among non-oracle variants. An exact-distance variant can be used only as an oracle upper bound because it relies on exact Euclidean distance information and is not a deployable approximate index.

\begin{table}[htbp]
\centering
\caption{Structure ablation of MECH approximate annulus query on ISOLET.}
\label{tab:mech_structure_ablation}
\resizebox{\linewidth}{!}{%
\begin{tabular}{lccccccc}
\toprule
Variant & P & R & $R^k$ & $E_+^k$ & KDE Err. & CER & Online ms \\
\midrule
Hash Only & 10.00\% & 100.00\% & 100.00\% & 750.02\% & 750.02\% & 10.00 & 0.849 \\
Hash + Sphere & 0.00\% & 0.00\% & 0.00\% & 0.00\% & 100.00\% & 0.00 & 0.261 \\
Hash + Distance Table & 66.73\% & 72.35\% & 72.37\% & 35.38\% & 9.34\% & 1.09 & 0.553 \\
Full MECH & \textbf{96.89\%} & \textbf{96.75\%} & \textbf{96.76\%} & \textbf{3.10\%} & \textbf{1.43\%} & 1.00 & 0.660 \\
\bottomrule
\end{tabular}%
}
\end{table}

\subsection{MECH Parameter Sensitivity}
Table~\ref{tab:mech_sensitivity} reports single-factor sensitivity for the main MECH pipeline. After origin-shift preprocessing, the effect of $L$ and $c$ becomes weak because the sphere-layer and angular filters already preserve most valid annulus points. The sphere interval $\delta$ remains important: a finer interval such as $\delta=0.25$ reduces radial quantization error and yields higher recall and lower KDE error. The annulus radius factor $\rho$ changes the queried distance region, while the bandwidth factor $\sigma$ changes only the kernel weights of the same candidate set.

\begin{table}[htbp]
\centering
\caption{MECH parameter sensitivity with kernel-weighted KDE metrics on ISOLET.}
\label{tab:mech_sensitivity}
\resizebox{\linewidth}{!}{%
\begin{tabular}{lcccccc}
\toprule
Setting & P & R & $R^k$ & KDE Err. & CER & Online ms \\
\midrule
$L=2$ & 96.89\% & 96.75\% & 96.76\% & 1.43\% & 1.00 & 0.656 \\
$L=4$ & 96.89\% & 96.75\% & 96.76\% & 1.43\% & 1.00 & 0.651 \\
$L=8$ & 96.89\% & 96.75\% & 96.76\% & 1.43\% & 1.00 & 0.677 \\
$L=12$ & 96.89\% & 96.75\% & 96.76\% & 1.43\% & 1.00 & 0.960 \\
$c=4$ & 96.89\% & 96.75\% & 96.76\% & 1.43\% & 1.00 & 0.676 \\
$c=8$ & 96.89\% & 96.75\% & 96.76\% & 1.43\% & 1.00 & 0.655 \\
$c=12$ & 96.89\% & 96.75\% & 96.76\% & 1.43\% & 1.00 & 0.653 \\
$\delta=0.25$ & 96.89\% & 96.75\% & 96.76\% & 1.43\% & 1.00 & 0.662 \\
$\delta=0.5$ & 94.59\% & 93.83\% & 93.85\% & 1.89\% & 0.99 & 0.639 \\
$\delta=1.0$ & 89.08\% & 88.17\% & 88.20\% & 3.06\% & 0.99 & 0.650 \\
$\delta=2.0$ & 77.26\% & 76.10\% & 76.19\% & 4.63\% & 0.99 & 0.642 \\
$\rho=0.75$ & 97.00\% & 97.79\% & 97.78\% & 3.60\% & 1.01 & 0.560 \\
$\rho=1.0$ & 96.89\% & 96.75\% & 96.76\% & 1.43\% & 1.00 & 0.642 \\
$\rho=1.25$ & 97.38\% & 97.55\% & 97.56\% & 1.30\% & 1.00 & 0.745 \\
$\rho=1.5$ & 98.68\% & 97.52\% & 97.45\% & 1.86\% & 0.99 & 0.663 \\
$\sigma=0.5$ & 96.89\% & 96.75\% & 96.78\% & 1.43\% & 1.00 & 0.645 \\
$\sigma=1.0$ & 96.89\% & 96.75\% & 96.76\% & 1.43\% & 1.00 & 0.644 \\
$\sigma=2.0$ & 96.89\% & 96.75\% & 96.75\% & 1.43\% & 1.00 & 0.657 \\
\bottomrule
\end{tabular}%
}
\end{table}

\subsection{Kernel Density Estimation for Anomaly Detection}
This study conducts experimental validation using the real-world UNSW-NB15 network traffic dataset. The dataset contains 49 features and distinguishes normal traffic from attack traffic. A total of 10,000 samples are randomly selected. Non-numerical features are removed, and Z-score normalization is applied. The algorithm employs approximate annulus partition combined with Gaussian KDE, and the decision threshold is selected by grid search on the validation set.

Experimental results in Fig.~\ref{fig:unsw1} show that predictions from both accurate annulus and approximate annulus methods align closely with the ground truth, whereas MCMC exhibits significant deviations. Table~\ref{tab:performance} shows that approximate annulus KDE achieves the best overall anomaly detection performance, with F1=98.1\%, Precision=96.2\%, and Recall=100\%, while being 15.6 times faster than accurate annulus KDE.

\begin{figure}
    \centering
    \includegraphics[width=1\linewidth]{image/UNSW1.png}
    \caption{t-SNE visualization of UNSW-NB15: (a) ground truth, (b) accurate annulus predictions, (c) approximate annulus predictions, and (d) MCMC predictions.}
    \label{fig:unsw1}
\end{figure}

\begin{table}[htbp]
\centering
\caption{Comparative performance evaluation of anomaly detection methods on UNSW-NB15.}
\resizebox{\linewidth}{!}{%
\begin{tabular}{lccccc}
\toprule
Methods & Threshold & P & R & F1 & Time \\
\midrule
Accurate Annulus    & 5    & 93.0\% & 100\% & 96.4\% & 13.6s \\
Approximate Annulus & 5    & \textbf{96.2\%} & \textbf{100\%} & \textbf{98.1\%} & \textbf{0.87s} \\
MCMC                & 0.02 & 81.5\% & 8.6\%  & 15.7\% & 10.1s \\
\bottomrule
\end{tabular}%
}
\label{tab:performance}
\end{table}
\section{Conclusion and Future Work}


This paper addresses the problem of kernel density estimation for high-dimensional and large-scale data in the field of consumer electronics and proposes an efficient estimation framework based on approximate annulus partition. Under a given error constraint, the method effectively aggregates samples with similar radial distances, requiring only a small number of samples within each annulus to estimate the kernel density contribution, thereby significantly reducing the sampling scale and computational cost. Experimental results demonstrate that the proposed method achieves superior performance across multiple high-dimensional datasets and network traffic anomaly detection tasks.

Future work will focus on three main directions: optimizing the indexing and parameter mechanisms; establishing a systematic theoretical analysis framework to quantitatively characterize error propagation and computational complexity; designing data-dependent adaptive partition strategies that jointly optimize annulus parameters and kernel bandwidth. These extensions are expected to further enhance the applicability and impact of the proposed method in intelligent terminal data modeling and anomaly detection applications within the consumer electronics domain.


 % 或IEEEtran  elsarticle-num, plain, gbt7714-numerical 
\bibliographystyle{IEEEtran}  
\bibliography{cas-refs} 

\end{document}
